\documentclass[10pt,a4paper]{amsart}
\usepackage[margin=19mm]{geometry}
\usepackage{amsmath,amssymb,mathtools}
\usepackage[scr=rsfs,cal=euler]{mathalfa}
\usepackage{xfrac}
\usepackage[unicode,hidelinks,bookmarks=false]{hyperref}
\newcommand{\ie}{i.e.\ }
\newcommand{\cf}{cf.\ }
\newcommand{\resp}{respectively\ }
\newcommand{\Subsection}{\subsection}
\providecommand{\nobreakdash}{\nobreak\hbox{-}}
\setlength{\emergencystretch}{2em}
\multlinegap0pt
\usepackage{amssymb}
\usepackage{mathtools}
\usepackage{tikz}
\usepackage{xcolor}
\newtheorem{theorem}{Theorem}
\title{Exact expressions for the unresolved stress in a finite-volume based large-eddy simulation}
\author{Syver D{\o}ving Agdestein, Roel Verstappen, Benjamin Sanderse}
\date{Source: arXiv:2507.17051v2 (CC BY 4.0)}
\begin{document}
\maketitle

\noindent\emph{Source and licence.} Exact expressions for the unresolved stress in a finite-volume based large-eddy simulation. Version arXiv:2507.17051v2 (CC BY 4.0), distributed by arXiv under the Creative Commons Attribution 4.0 International licence, http://creativecommons.org/licenses/by/4.0/, as recorded in the arXiv licence metadata for this exact version and shown on the abstract page https://arxiv.org/abs/2507.17051v2. Full licence notice: \texttt{licenses/arxiv-2507.17051-CC-BY-4.0.txt}.\par\smallskip

\noindent\textit{Editorial scope.} Counted unit: the article's theorem th:stress-projector (source0.tex lines 3259--3265). The article's complete original proof of it is delivered with it (source0.tex lines 3267--3321, one \texttt{proof} environment of 55 lines). No mathematics is written, repaired or re-derived: the statement and the proof are the article's own lines, in the article's own notation, and no retained line is substituted or re-flowed.  No further source-local context is needed: every cross-reference of every retained line resolves inside the two delivered spans.  Nothing was omitted from any retained span, so the package carries no omission record.  No retained line contains a citation, so the wrapper carries no bibliography and no bibliography entry.  No retained line contains a figure environment, an image-inclusion command or drawing code, so no image is delivered and the package declares no asset dependency.  Editorial formatting only, in addition: an AMS \texttt{amsart} preamble that restates the article's own declarations and macros used by the retained lines, namely the environments \texttt{theorem} on separate counters, together with the standard packages those lines need.  The retained environments print in this document as Theorem 1 in order of appearance, whereas the article prints the same items with the numbering of its own full document.  The complete native source of the article as distributed in the arXiv e-print for this exact version is bundled unchanged as \texttt{sources/182052/source0.tex}.
\begin{theorem} \label{th:stress-projector}
    The operator $\pi_{i j \alpha \beta}$ is a projector onto the space of
    divergence-preserving stress tensors, i.e.\
    $\partial_i \partial_j \pi_{i j \alpha \beta} = 0$
    (the output of $\pi$ is divergence-preserving) and
    $\pi \pi = \pi$.
\end{theorem}

\begin{proof}
    The double-divergence $\partial_i \partial_j$ composed with the operator
    $\pi_{i j \alpha \beta}$ is
    \begin{equation}
    \begin{split}
        \partial_i \partial_j \pi_{i j \alpha \beta}
        & = \delta_{i \alpha} \delta_{j \beta} \partial_i \partial_j
        - \delta_{i j} \partial_i \partial_j
        \left( \partial_k \partial_k \right)^{\dagger}
        \partial_\alpha \partial_\beta \\
        & = \partial_\alpha \partial_\beta
        - \underbrace{\partial_i \partial_i
        \left( \partial_k \partial_k \right)^{\dagger}}_{1}
        \partial_\alpha \partial_\beta \\
        & = \partial_\alpha \partial_\beta - \partial_\alpha \partial_\beta \\
        & = 0.
    \end{split}
    \end{equation}
    This means that for all $\sigma \in U^{3 \times 3}$, we have
    $\partial_i \partial_j ( \pi_{i j \alpha \beta} \sigma_{\alpha \beta} ) = 0$,
    so $\pi \sigma$ is a divergence-preserving
    tensor.

    Furthermore, applying the operator twice gives
    \begin{equation}
    \begin{split}
        \pi_{i j \alpha \beta} \pi_{\alpha \beta m n} =
        & \left( \delta_{i \alpha} \delta_{j \beta} - \delta_{i j} (\partial_k \partial_k)^{\dagger}
        \partial_\alpha \partial_\beta \right) \\
        & \left( \delta_{\alpha m} \delta_{\beta n} - \delta_{\alpha \beta} (\partial_l \partial_l)^{\dagger} \partial_m \partial_n \right) \\
        = &
        (\delta_{i \alpha} \delta_{j \beta}) (\delta_{\alpha m} \delta_{\beta n}) \\
        - &
        (\delta_{i \alpha} \delta_{j \beta}) \delta_{\alpha \beta} (\partial_l \partial_l)^{\dagger} \partial_m \partial_n \\
        - &
        (\delta_{\alpha m} \delta_{\beta n}) \delta_{i j} (\partial_k \partial_k)^{\dagger} \partial_\alpha \partial_\beta \\
        + &
        \delta_{i j} \delta_{\alpha \beta}
        (\partial_k \partial_k)^{\dagger} \partial_\alpha \partial_\beta
        (\partial_l \partial_l)^{\dagger} \partial_m \partial_n \\
        = &
        \delta_{i m} \delta_{j n} \\
        - & \delta_{i j} (\partial_l \partial_l)^{\dagger} \partial_m \partial_n \\
        - & \delta_{i j} (\partial_k \partial_k)^{\dagger} \partial_m \partial_n \\
        + & \delta_{i j} (\partial_k \partial_k)^{\dagger}
        \underbrace{\partial_\alpha \partial_\alpha (\partial_l
        \partial_l)^{\dagger}}_{1}
        \partial_m \partial_n \\
        = & \delta_{i m} \delta_{j n} + (-2 + 1) \delta_{i j} (\partial_k \partial_k)^{\dagger} \partial_m \partial_n \\
        = & \pi_{i j m n}.
    \end{split}
    \end{equation}
    Since $\pi \pi = \pi$, we
    can conclude that $\pi$ is idempotent.
\end{proof}

\end{document}
